\section{Theoretical Motivation}
\label{sec:theoretical-motivation}

\subsection{Governance as an emergent relational ecology}
\label{sec:emergent-ecology}

Agent-based computational social science treats macro-level social pattern
as an emergent consequence of micro-level rules, rather than as a
directly-specified target
\parencite{epstein1996growing,railsback2019agent}. The tradition most
directly relevant here descends from two distinct lines: iterated-game
accounts of cooperation, reciprocity, and costly punishment
\parencite{axelrod1984evolution,axelrod1981evolution,fehr2002altruistic},
and empirically grounded common-pool-resource governance research
\parencite{ostrom1990governing,janssen2006empirically,poteete2010working}.
Ostrom's fieldwork identified design principles --- clearly defined
boundaries, congruence between rules and local conditions, collective-choice
arrangements, monitoring, graduated sanctions, conflict-resolution
mechanisms, and nested enterprises --- under which self-governing
institutions for shared resources persist without either central
authority or privatization. Our model does not encode these principles
directly; it encodes the more elementary primitives (extraction,
maintenance, monitoring, sanctioning, delegation, coordination) from which
such principles, if they matter, would have to emerge as statistical
regularities in simulated outcomes, not as assumed structure.

\subsection{From atomic interactions to proto-institutions}
\label{sec:atomic-to-proto}

The project's own development sequence (\cref{sec:model-genealogy}) is
built explicitly bottom-up: a population of heterogeneous agents interacts
bilaterally under a repeated cooperation/defection grammar with imperfect
observation, memory, and optional costly punishment; the population is
coupled to a materially accounted resource ecology; and only once that
ecology's own dynamics were understood did the project introduce a
locally shared, voluntarily adopted, imperfectly monitored and enforced
rule --- a \textit{proto-institution} --- as an opt-in generative
mechanism. Proto-institutions are then extended with delegation
(administrative capacity accumulated from realized performance, never
granted as a trait), discretion and its accountability counterweight
(audit, appeal, recall), and interpatch coordination. At every stage, the
mechanism defaults to a verified, byte-for-byte no-op, so that each new
layer is additive rather than a silent redefinition of what came before.

\subsection{Recurrent institutional phenotypes versus imposed typologies}
\label{sec:recurrent-vs-imposed}

A central methodological commitment of this project, enforced structurally
rather than merely stated, is that the generative layer never encodes a
governance typology, and the observer layer never imposes one on the
results before the results are seen. We use \textit{recurrent
institutional phenotype} throughout, rather than \textit{institutional
type} or \textit{governance regime}, to mark this distinction: a phenotype
is a statistically discovered pattern in an observer-computed feature
vector, evaluated against explicit stability and recurrence criteria; it is
never a category the model was built to produce. This distinguishes our
approach from ABM applications that calibrate agent rules to reproduce a
known macro-target. Here the target category count is itself an open
empirical question the discovery pipeline is asked, not told.

\subsection{Historical institutionalism, path dependence, and critical junctures}
\label{sec:historical-institutionalism}

A separate literature, largely orthogonal to the agent-based tradition
above, argues that governance outcomes are shaped not only by
contemporaneous structural conditions but by the specific historical
sequence through which institutions formed: critical junctures at which
a range of outcomes was genuinely possible, followed by increasing-returns
dynamics that make an early institutional choice progressively more costly
to reverse \parencite{pierson2000increasing,mahoney2000path,capoccia2007critical};
the long shadow of colonial extraction versus
settlement patterns and early state-building choices on later institutional
quality \parencite{acemoglu2012why,northwallisweingast2009}; and the
broader claim that political order and its erosion are themselves
historically contingent processes, not fixed endpoints
\parencite{fukuyama2014political,north1990institutions}. This literature
supplies the theoretical vocabulary --- sequencing, entrenchment, critical
juncture, lock-in --- that Level 1.3/1.3b's own historical mechanisms
operationalize generatively (\cref{sec:level13}); it does not, by itself,
adjudicate whether path-dependent processes compress governance outcomes
into a small number of discrete, recurring configurations or instead
produce a continuously varying space with graded, historically contingent
structure. That is an empirical question this paper's own historical
discovery program is built to ask without presupposing the answer, in the
same outcome-blind spirit as Level 1.2/1.2b's synchronic search.

The dissertation's own governance-indicator critique is relevant here too.
Aggregate governance rankings such as the World Bank's Worldwide Governance
Indicators \parencite{kaufmann2011worldwide} compress multiple correlated
component indicators onto a single latent dimension via an unobserved
components model, an approach the dissertation argues collapses
theoretically distinct dimensions --- capacity to act and constraint on
how that capacity may be used --- into one ranking, and which is, by
construction, silent on historical trajectory: a single cross-sectional
score cannot represent how a state arrived at its current position, only
where it currently sits. The dissertation's own two-dimensional $S$/$Z$
reframing is explicitly a response to this limitation; this paper's own
historical extension is a further response, adding a third
axis --- historical trajectory itself --- that neither a one-dimensional
ranking nor a two-dimensional cross-sectional coordinate space can
represent.

\subsection{Terminological discipline}
\label{sec:terminology-discipline}

This paper uses five terms that name related but analytically distinct
objects, and treats confusing them as a substantive error rather than a
stylistic one, given how directly this project's own headline findings turn
on the distinction. A \textit{configuration} is a specific, realized point
in the observer-computed feature space (Level 1.2/1.2b's ten-dimensional
phenotype vector, or Level 1.3/1.3b's fourteen-dimensional extension) ---
one simulated jurisdiction's own outcome, nothing more. A
\textit{statistical component} is a fitted element of a clustering or
mixture model (a Gaussian-mixture component, a hierarchical-clustering
branch) --- a property of a fitting procedure applied to a finite sample,
not by itself evidence of anything in the underlying process. A
\textit{recurrent institutional phenotype} is a statistical component that
has additionally cleared this project's own preregistered recurrence
criteria (\cref{sec:discovery-procedure}) --- cross-seed recurrence,
cross-algorithm support, bootstrap stability, and the rest --- evaluated
by strict conjunction, never a majority. A \textit{dynamical attractor} is
a further, stronger claim this project has never made about any discovered
structure: that a system returns to a specific state after perturbation,
which requires direct, disclosed evidence of majority recovery on every
tracked dimension (\cref{sec:level12b-scale-causal}), not merely that a
state persists once reached. An \textit{empirical trajectory}, finally, is
the dissertation's own object: a qualitatively interpreted region of its
SEM-identified $S$/$Z$ space, populated by real countries' survey-measured
positions, which this paper's own simulation results can inform but cannot
directly confirm or refute (\cref{sec:disc-relationship-to-dissertation}).

We use \textit{historically structured, heterogeneous governance
continuum} as this paper's own preferred summary phrase for what Level
1.3/1.3b's evidence supports, and, where the evidence licenses a
marginally stronger reading, \textit{a continuous institutional space
containing conditional path dependence and possible metastable regions}
--- never stronger than that. We do not, anywhere in this paper, claim that
the model discovered exactly three or exactly eight governance species,
that the ABM's own discovery pipeline found three attractors corresponding
to the dissertation's trajectories, that path dependence in this model is
universal or historically deterministic, or that any simulation result in
this paper empirically validates the dissertation's own trajectory
typology. Each of these five specific claims is one this paper's own
evidence does not support, and \cref{sec:disc-boundary-claims} states
explicitly, alongside the positive findings, why each is ruled out.

\subsection{Why null discovery is theoretically informative}
\label{sec:why-null-informative}

A discovery procedure that always finds structure is not informative about
whether structure exists; it is informative only about the procedure's own
bias. The clustering-methodology literature treats this directly: cluster
stability under resampling and under alternative algorithms is itself the
primary evidence for whether a partition reflects real structure or
artifact of the fitting procedure
\parencite{benhur2002stability,vonluxburg2010clustering,hubert1985comparing,rousseeuw1987silhouettes}.
Model-based (finite-mixture) clustering in particular is well known to
overfit at modest sample sizes, producing spuriously confident partitions
that do not generalize to held-out data \parencite{fraley2002model};
density-based clustering supplies a structurally different failure mode
(explicit noise, no forced full-coverage partition) precisely so that a
finding of "structure" is not an artifact of a single clustering
family's own inductive bias \parencite{campello2013density}. Reporting a
well-designed null result --- one that would have detected recurrent
structure had it been there, evaluated by multiple algorithms and
independently replicated on fresh data --- is therefore not merely
permissible negative-results hygiene; in a discovery setting it is direct,
positive evidence about the shape of the underlying outcome space. This is
the logic \cref{sec:level12,sec:level12b} apply throughout: every reported
null is a preregistered, multiply-corroborated, independently replicated
finding, not an absence of search.
